\documentclass[journal]{IEEEtran}

% -----------------------------------------------------------------------
% Packages
% -----------------------------------------------------------------------
\usepackage{amsmath}
\usepackage{amssymb}
\usepackage{graphicx}
\usepackage{cite}
\usepackage{hyperref}
\usepackage{booktabs}
\usepackage{algorithm}
\usepackage{algpseudocode}
\usepackage{bm}
\usepackage{xcolor}
\usepackage{subcaption}
\usepackage{multirow}
\usepackage{array}
\usepackage{url}

% -----------------------------------------------------------------------
% Document
% -----------------------------------------------------------------------
\begin{document}

\title{Deep Neural Network Surrogate Model for Multi-Period AC Optimal Power Flow}

\author{%
  \IEEEauthorblockN{Author Name}\\
  \IEEEauthorblockA{Department of Electrical Engineering\\
  University Name, City, Country\\
  email@university.edu}%
}

\maketitle

% -----------------------------------------------------------------------
% Abstract
% -----------------------------------------------------------------------
\begin{abstract}
This paper presents a surrogate-based optimization framework for the Multi-Period
AC Optimal Power Flow (MP-ACOPF) problem. A deep neural network (DNN) is trained
to replace the computationally expensive AC power flow equations, enabling faster
optimization while maintaining physical consistency. The trained surrogate is
embedded into a nonlinear optimization framework solved via the Augmented Lagrangian
Method (ALM), which penalizes power balance constraint violations and guides the
optimizer towards physically feasible solutions. Accurate prediction of the power
flow variables — particularly voltage magnitudes and phase angles at the buses — is
identified as a critical requirement, since even small prediction errors can amplify
into large power balance mismatches due to the inherent nonlinearity of the AC power
flow equations. Numerical experiments on benchmark IEEE power systems demonstrate
the effectiveness of the proposed approach in terms of convergence, constraint
satisfaction, and solution quality compared to direct nonlinear programming solvers.
\end{abstract}

\begin{IEEEkeywords}
AC optimal power flow, deep neural network, surrogate model, multi-period
optimization, augmented Lagrangian method, voltage magnitude, power balance.
\end{IEEEkeywords}

% -----------------------------------------------------------------------
% I. Introduction
% -----------------------------------------------------------------------
\section{Introduction}
\label{sec:introduction}

The optimal power flow (OPF) problem is a cornerstone of power system operation
and planning. It determines the most economical dispatch of controllable generation
units while satisfying a set of physical, network, and operational constraints. The
AC formulation (ACOPF) models the nonlinear power flow physics faithfully, making
it the most accurate variant, but also the most computationally demanding due to
its nonconvex, nonlinear structure~\cite{carpentier1962contribution,frank2012optimal}.

The Multi-Period AC Optimal Power Flow (MP-ACOPF) extends the classical single-period
ACOPF by optimizing over multiple consecutive time periods simultaneously. This
extension is essential for capturing temporal dependencies that are increasingly
important in modern power systems: generator ramping limits, energy storage charge
and discharge cycles, and the variability of renewable energy sources such as wind
and solar. By solving over the full time horizon at once, MP-ACOPF enables globally
coherent schedules that single-period methods cannot guarantee. However, the
simultaneous treatment of all time periods multiplies the problem size by a factor
of $T$ (the number of periods) and introduces inter-period coupling constraints,
dramatically increasing the computational burden~\cite{jabr2012adjustable}.
Solving MP-ACOPF reliably and efficiently is therefore critical for day-ahead
scheduling, real-time grid management, and the economic integration of
renewable generation.

Classical approaches to ACOPF include interior-point methods, sequential quadratic
programming, and successive linear programming, all of which can struggle with large
problem instances or tight time requirements~\cite{capitanescu2011interior}.
Convex relaxation techniques, notably semidefinite programming (SDP) and
second-order cone programming (SOCP) relaxations, provide tractable approximations
but may return inexact or infeasible solutions for general meshed
networks~\cite{lavaei2012zero}. The DC power flow linearization sacrifices reactive
power modeling and voltage accuracy in exchange for computational simplicity, and
is therefore unsuitable when precise voltage profiles are required~\cite{stott2009dc}.
These trade-offs motivate the search for new computational paradigms that can
deliver high-quality ACOPF solutions at reduced cost.

Machine learning — and in particular deep neural networks (DNNs) — has emerged as
a promising direction for accelerating power system computations. DNNs can learn
complex nonlinear input-output mappings from data, and once trained, they provide
virtually instantaneous inference. Several recent works exploit DNNs to predict
near-optimal OPF solutions directly from load
inputs~\cite{zamzam2020learning,pan2023deepopf}, to approximate constraint
functions~\cite{venzke2020learning}, or to warm-start nonlinear solvers. The
concept of \emph{DNN surrogate models} — replacing expensive simulation or
constraint-evaluation steps with a trained network — has gained significant
traction as a way to embed learned physical representations into optimization
frameworks~\cite{chen2020learning,donti2021dc3}.

In the context of surrogate-based OPF, the DNN replaces the full AC power flow
equations within the optimization loop, allowing gradient-based solvers to operate
on a smooth, differentiable approximation of the power flow constraints. This
yields substantial speed-ups compared to repeatedly invoking a full power flow
solver. The key challenge is ensuring that the surrogate captures the relevant
physics with sufficient accuracy so that the optimizer converges to a solution
that is also feasible for the true AC power flow. This becomes particularly
demanding in the multi-period setting, where prediction errors can accumulate
across time steps and compound through the coupling constraints.

Several recent contributions have demonstrated the viability of DNN surrogate
models for power system optimization. Zamzam and Baker~\cite{zamzam2020learning}
showed that DNNs can predict near-optimal ACOPF solutions with high accuracy.
Venzke~et~al.~\cite{venzke2020learning} combined DNN surrogates with probabilistic
safety certificates. Donti~et~al.~\cite{donti2021dc3} proposed DC3, a
differentiable framework that enforces constraint satisfaction through learned
completion layers. Pan~et~al.~\cite{pan2023deepopf} presented DeepOPF, scaling
DNN-based OPF prediction to large systems while maintaining strong generalization.
Despite these advances, the extension to MP-ACOPF and the detailed analysis of
how DNN prediction accuracy affects optimization convergence remain active areas
of research.

A critical aspect of surrogate-based MP-ACOPF is accurately predicting the power
flow variables — in particular, the voltage magnitude and phase angle at each bus
— since even small prediction errors can lead to high power balance mismatches due
to the nonlinearity of the AC power flow equations.

The proposed framework addresses these challenges by training a DNN on power flow
simulation data generated via \texttt{pandapower}~\cite{thurner2018pandapower},
and embedding the trained surrogate into an Augmented Lagrangian Method (ALM)
optimization loop. The ALM provides a natural mechanism to penalize constraint
violations, and the DNN surrogate enables efficient gradient computation through
automatic differentiation. Initialization of the optimization from a feasible
power flow solution further improves convergence and reduces initial equality
constraint violations.

The remainder of this paper is organized as follows.
Section~\ref{sec:problem} formulates the MP-ACOPF problem.
Section~\ref{sec:surrogate} describes the DNN surrogate model and training.
Section~\ref{sec:optimization} presents the surrogate-based optimization framework.
Section~\ref{sec:experiments} reports numerical experiments, and
Section~\ref{sec:conclusion} concludes.

% -----------------------------------------------------------------------
% II. Problem Formulation
% -----------------------------------------------------------------------
\section{Problem Formulation}
\label{sec:problem}

The MP-ACOPF problem over a planning horizon of $T$ time periods is formulated as:
\begin{align}
  \min_{\mathbf{x}_1, \ldots, \mathbf{x}_T} \quad
    & \sum_{t=1}^{T} f_t(\mathbf{x}_t) \label{eq:obj} \\
  \text{s.t.} \quad
    & g_t(\mathbf{x}_t) = 0, \quad t = 1, \ldots, T \label{eq:eq} \\
    & h_t(\mathbf{x}_t) \leq 0, \quad t = 1, \ldots, T \label{eq:ineq} \\
    & c(\mathbf{x}_t, \mathbf{x}_{t+1}) \leq 0, \quad t = 1, \ldots, T{-}1
      \label{eq:coupling}
\end{align}

where $\mathbf{x}_t$ collects the decision variables at time $t$ (active and
reactive power injections, voltage magnitudes, and phase angles), $f_t$ denotes
the generation cost, $g_t$ represents the power balance equality constraints,
$h_t$ includes operational inequality constraints (voltage limits, line flow limits,
generator capacity limits), and $c$ captures inter-period coupling constraints
such as ramping limits.

\subsection{AC Power Flow Equations}

The AC power flow equations for bus $i$ at time $t$ read:
\begin{align}
  P_{i,t} &= V_{i,t} \sum_{j \in \mathcal{N}}
    V_{j,t} \left( G_{ij} \cos\theta_{ij,t} + B_{ij} \sin\theta_{ij,t} \right)
    \label{eq:pf_p} \\
  Q_{i,t} &= V_{i,t} \sum_{j \in \mathcal{N}}
    V_{j,t} \left( G_{ij} \sin\theta_{ij,t} - B_{ij} \cos\theta_{ij,t} \right)
    \label{eq:pf_q}
\end{align}

where $P_{i,t}$ and $Q_{i,t}$ are the net active and reactive power injections,
$V_{i,t}$ is the voltage magnitude, $\theta_{ij,t} = \theta_{i,t} - \theta_{j,t}$
is the voltage angle difference, and $G_{ij}$, $B_{ij}$ are the conductance and
susceptance entries of the nodal admittance matrix.

% -----------------------------------------------------------------------
% III. DNN Surrogate Model
% -----------------------------------------------------------------------
\section{DNN Surrogate Model}
\label{sec:surrogate}

\subsection{Dataset Generation}

A dataset of power flow solutions is generated by sampling load profiles from a
predefined distribution and solving the corresponding power flow using
\texttt{pandapower}. Each sample consists of the input vector $\mathbf{u}_t$
(active and reactive load demands) and the output vector $\mathbf{y}_t$ (voltage
magnitudes and phase angles at all buses, and generator active and reactive power).
The dataset spans a representative range of operating conditions to ensure
surrogate generalization.

\subsection{Network Architecture and Training}

A fully connected DNN with residual connections is employed. The network maps
load inputs to power flow state variables. Multiple hidden layers with ReLU
activations are trained by minimizing the mean squared error (MSE) on the dataset
using the Adam optimizer. Batch normalization and dropout regularize training
and prevent overfitting.

% -----------------------------------------------------------------------
% IV. Surrogate-Based Optimization
% -----------------------------------------------------------------------
\section{Surrogate-Based Optimization}
\label{sec:optimization}

\subsection{Augmented Lagrangian Method}

The DNN surrogate replaces~\eqref{eq:pf_p}--\eqref{eq:pf_q} in the optimization
loop. The resulting surrogate-based problem is solved via the ALM:
\begin{align}
  \mathcal{L}_\rho(\mathbf{x}, \bm{\lambda})
    = \sum_{t=1}^{T} f_t(\mathbf{x}_t)
    + \bm{\lambda}^\top \hat{g}(\mathbf{x})
    + \frac{\rho}{2} \|\hat{g}(\mathbf{x})\|^2
\end{align}

where $\hat{g}(\mathbf{x})$ denotes the surrogate-approximated power balance
residuals, $\bm{\lambda}$ are the Lagrange multipliers, and $\rho > 0$ is the
penalty parameter updated across ALM outer iterations.

\subsection{Initialization}

The optimization is initialized from a feasible power flow solution for the nominal
load, obtained from \texttt{pandapower}. This warm-start significantly reduces
initial equality constraint violations and improves ALM convergence.

% -----------------------------------------------------------------------
% V. Numerical Experiments
% -----------------------------------------------------------------------
\section{Numerical Experiments}
\label{sec:experiments}

Experiments are conducted on standard IEEE bus test systems. The DNN surrogate
is evaluated in terms of prediction accuracy and the quality of the resulting
optimal solutions relative to a direct nonlinear programming solver.

% -----------------------------------------------------------------------
% VI. Conclusion
% -----------------------------------------------------------------------
\section{Conclusion}
\label{sec:conclusion}

This paper presented a DNN surrogate-based framework for MP-ACOPF, embedding a
trained neural network into an ALM optimization loop. Accurate prediction of
voltage magnitudes and phase angles is essential: even small errors are amplified
by the nonlinear power flow equations into large power balance mismatches. Careful
initialization and high surrogate fidelity are identified as key factors for
convergence. Future work will investigate adaptive surrogate updating and
extensions to larger systems.

\bibliographystyle{IEEEtran}
\bibliography{references}

\end{document}
